\documentclass[12pt,a4paper]{report}

%% Vložení užitečných balíčků
%%% Hlavní soubor. Zde se definují základní parametry a~odkazuje se na ostatní části. %%%

%% Verze pro jednostranný tisk:
% Okraje: levý 40mm, pravý 25mm, horní a~dolní 25mm
% (ale pozor, LaTeX si sám přidává 1in)
\setlength\textwidth{145mm}
\setlength\textheight{247mm}
\setlength\oddsidemargin{15mm}
\setlength\evensidemargin{15mm}
\setlength\topmargin{0mm}
\setlength\headsep{0mm}
\setlength\headheight{0mm}
% \openright zařídí, aby následující text začínal na pravé straně knihy
\let\openright=\clearpage

%% Pokud tiskneme oboustranně:
% \documentclass[12pt,a4paper,twoside,openright]{report}
% \setlength\textwidth{145mm}
% \setlength\textheight{247mm}
% \setlength\oddsidemargin{14.2mm}
% \setlength\evensidemargin{0mm}
% \setlength\topmargin{0mm}
% \setlength\headsep{0mm}
% \setlength\headheight{0mm}
% \let\openright=\cleardoublepage

\usepackage[pdfa,colorlinks=false,urlcolor=black,bookmarksopen=true]{hyperref} 	% Záložky v~souboru

%% Vytváříme PDF/A-2u
\usepackage[a-2u]{pdfx}

%% Přepneme na českou sazbu a~fonty Latin Modern
\usepackage[czech]{babel}
\usepackage{lmodern}
\usepackage[T1]{fontenc}
\usepackage{textcomp}

%% Použité kódování znaků: obvykle latin2, cp1250 nebo utf8:
\usepackage[utf8]{inputenc}
\usepackage{csquotes}
\MakeOuterQuote{"}

%%% Další užitečné balíčky (jsou součástí běžných distribucí LaTeXu)
\usepackage{amsmath}        % rozšíření pro sazbu matematiky
\usepackage{amsfonts}       % matematické fonty
\usepackage{amsthm}         % sazba vět, definic apod.
\usepackage{amssymb}
\usepackage{bbding}         % balíček s nejrůznějšími symboly
			    % (čtverečky, hvězdičky, tužtičky, nůžtičky, ...)
\usepackage{bm}             % tučné symboly (příkaz \bm)
\usepackage{graphicx}       % vkládání obrázků
\usepackage{fancyvrb}       % vylepšené prostředí pro strojové písmo
\usepackage{indentfirst}    % zavede odsazení 1. odstavce kapitoly
\usepackage[numbers]{natbib}         % zajištuje možnost odkazovat na literaturu
			    % stylem AUTOR (ROK), resp. AUTOR [ČÍSLO]
\usepackage[nottoc]{tocbibind} % zajistí přidání seznamu literatury,
                            % obrázků a~tabulek do obsahu
\usepackage{icomma}         % inteligetní čárka v~matematickém módu
\usepackage{dcolumn}        % lepší zarovnání sloupců v~tabulkách
\usepackage{booktabs}       % lepší vodorovné linky v~tabulkách
\usepackage{paralist}       % lepší enumerate a~itemize
\usepackage{xcolor}         % barevná sazba
\usepackage{float}
\usepackage{ifthen}
\usepackage{cancel}
\usepackage{enumitem}
\usepackage{ragged2e}
\usepackage{needspace}		
\usepackage{wrapfig}		% plovoucí obrázky nalevo/napravo
\usepackage{pgfplots}
\usepackage{subcaption}
\pgfplotsset{compat=1.15}
\usepackage{mathrsfs}
\usetikzlibrary{arrows}

\justifying
\graphicspath{{./components/images/}}

%% Definice různých užitečných maker (viz popis uvnitř souboru)

\DefineBibliographyStrings{czech}{ % Definování vlastního názvu pro bibliografii v češtině
  bibliography = Seznam použité literatury,
}

% TODO
\newcommand{\todo}[1]{\textcolor{red}{(\noindent TODO: #1.)}}

% Paths
\newcommand{\sectionpath}[1]{ch#1/sections}

% Celá jména
\newcommand{\name}[1]{\mbox{\textsc{#1}}}


% Cesty
\newcommand{\appendixpath}{components/appendix}

% Images
\newcommand{\subfigwidth}{6cm}        % šířka okénka pro "podobrázky"
\newcommand{\wrappedfigwidth}{6cm}
\newcommand{\fullhd}{0.2}             % měřítko Full HD obrázku
\newcommand{\normalipe}{0.8}          % standardní měřítko IPE obrázku
\newcommand{\fractalscale}{0.3}       % měřítko obrázku softwarem vygenerovaného fraktálu

%%% Prostředí pro sazbu kódu, případně vstupu/výstupu počítačových
%%% programů. (Vyžaduje balíček fancyvrb -- fancy verbatim.)

\DefineVerbatimEnvironment{code}{Verbatim}{fontsize=\small, frame=single}

% Zapne černé "slimáky" na koncích řádků, které přetekly, abychom si
% jich lépe všimli.
\overfullrule=1mm

% Trochu volnější nastavení dělení slov, než je default.
\lefthyphenmin=2
\righthyphenmin=2

\setlength{\parskip}{0.3em}



%%% Údaje o~práci

%%% Název práce v~jazyce práce (přesně podle zadání)
\def\NazevPrace{Proxmox}

%%% Název práce v~angličtině
\def\NazevPraceEN{\todo{Title}}

%%% Typ práce
\def\TypPrace{Seminární práce}

%%% Jméno autora
\def\AutorPrace{Martin Kopecký}



%%% Rok odevzdání
\def\RokOdevzdani{2024}

% Studijní program a~obor
\def\StudijniProgram{Aplikovaná informatika}
\def\StudijniObor{Informační systémy}


%%% 3 až 5 klíčových slov (doporučeno), každé uzavřeno ve složených závorkách
\def\KlicovaSlova{%
{}, {klicova slova}
}
%\def\KlicovaSlovaEN{%
%{keywords}
%}

%%% Balíček hyperref, kterým jdou vyrábět klikací odkazy v~PDF,
%%% ale hlavně ho používáme k~uložení metadat do PDF (včetně obsahu).
%%% Většinu nastavítek přednastaví balíček pdfx.
\hypersetup{unicode}
\hypersetup{breaklinks=true}


\renewcommand{\bibname}{Seznam použité literatury}
\addbibresource{./literature/literature.bib}

%% Titulní strana a~různé povinné informační strany
\begin{document}
%%% Titulní strana práce a další povinné informační strany

%%% Titulní strana práce

\pagestyle{empty}
\hypersetup{pageanchor=false}

\begin{center}

%\centerline{\mbox{\includegraphics[width=166mm]{components/images/ujep.png}}}
\null


{\bf\Large Univerzita Jana Evangelisty Purkyně \\ v Ústí nad Labem}

\vspace{2mm}

{\bf\Large Přírodovědecká fakulta}

\vspace{20mm}

\includegraphics[width=80mm]{ujep.png}

\vspace{20mm}


{\LARGE\bfseries\NazevPrace}

{\large\bfseries\TypPrace}

\vfill


\begin{flushleft}
\begin{tabular}{>{\bfseries}r l}
    \textbf{Vypracoval:} & \AutorPrace \\
    % Vedoucí práce:\Vedouci 
    % Konzultant: \Vedouci
    \\

    \textbf{Studijní program:}  &  \StudijniProgram \\
    \textbf{Studijní obor:}  &   \StudijniObor \\
\end{tabular}
\end{flushleft}

\vfill

% Zde doplňte rok
\large{Ú\raisebox{-0.85ex}{\textsuperscript{STÍ NAD}} L\raisebox{-0.85ex}{\textsuperscript{ABEM}} \RokOdevzdani}

\end{center}

%\newpage

%%% Následuje vevázaný list -- kopie podepsaného "Zadání diplomové práce".
%%% Toto zadání NENÍ součástí elektronické verze práce, nescanovat.

%%% Strana s čestným prohlášením k diplomové práci

%\openright
%\hypersetup{pageanchor=true}
%\pagestyle{plain}
%\pagenumbering{roman}
%\vglue 0pt plus 1fill

%\noindent
%Prohlašuji, že jsem tuto diplomovou práci vypracoval(a) samostatně a výhradně
%s~použitím citovaných pramenů, literatury a dalších odborných zdrojů.
%Tato práce nebyla využita k získání jiného nebo stejného titulu.
%
%\medskip\noindent
%Beru na~vědomí, že se na moji práci vztahují práva a povinnosti vyplývající
%ze zákona č. 121/2000 Sb., autorského zákona v~platném znění, zejména skutečnost,
%že Univerzita Karlova má právo na~uzavření licenční smlouvy o~užití této
%práce jako školního díla podle §60 odst. 1 autorského zákona.
%
%\vspace{10mm}
%
%\hbox{\hbox to 0.5\hsize{%
%V \hbox to 6em{\dotfill} dne \hbox to 6em{\dotfill}
%\hss}\hbox to 0.5\hsize{\dotfill\quad}}
%\smallskip
%\hbox{\hbox to 0.5\hsize{}\hbox to 0.5\hsize{\hfil Podpis autora\hfil}}
%
%\vspace{20mm}
%\newpage

%%% Poděkování

%\openright
%
%\noindent
%\Podekovani
%
%\newpage

%%% Povinná informační strana diplomové práce

%\openright
%
%\vbox to 0.5\vsize{
%\setlength\parindent{0mm}
%\setlength\parskip{5mm}
%
%Název práce:
%\NazevPrace
%
%Autor:
%\AutorPrace
%
%\TypPracoviste:
%\Katedra
%
%Vedoucí diplomové práce:
%\Vedouci, \KatedraVedouciho
%
%Abstrakt:
%\Abstrakt
%
%Klíčová slova:
%\KlicovaSlova
%
%\vss}\nobreak\vbox to 0.49\vsize{
%\setlength\parindent{0mm}
%\setlength\parskip{5mm}
%
%Title:
%\NazevPraceEN
%
%Author:
%\AutorPrace
%
%\TypPracovisteEN:
%\KatedraEN
%
%Supervisor:
%\Vedouci, \KatedraVedoucihoEN
%
%Abstract:
%\AbstraktEN
%
%Keywords:
%\KlicovaSlovaEN
%
%\vss}

\newpage

\openright
\pagestyle{plain}
\pagenumbering{arabic}
\setcounter{page}{1}

%%% Strana s automaticky generovaným obsahem bakalářské práce

\tableofcontents

%%% Jednotlivé kapitoly práce jsou pro přehlednost uloženy v~samostatných souborech
\chapter{Historie}

Proxmox je oblíbená virtualizační platforma, která kombinuje virtuální servery (Proxmox Virtual Environment, PVE) a systém pro správu clusterů a virtualizovaného úložiště (Proxmox Virtual Environment Cluster). Její historie sahá do roku 2008, kdy byla poprvé uvedena na trh společností Proxmox Server Solutions GmbH, založenou v Rakousku.

Proxmox VE byl vytvořen jako open-source řešení, které poskytuje efektivní a snadno použitelné prostředí pro správu virtuálních serverů. Od začátku byl postaven na Linuxovém jádru a využíval virtualizační technologie jako KVM (Kernel-based Virtual Machine) a kontejnerovou technologii LXC (Linux Containers), což umožňovalo uživatelům vytvářet a spravovat virtuální stroje a kontejnery z jediného rozhraní.

V průběhu let došlo k několika důležitým vývojovým krokům. Významným milníkem bylo zavedení Proxmox VE Cluster, což umožnilo správu více serverů jako jednoho celku, což výrazně zjednodušilo správu zdrojů a redundanci. Další významnou aktualizací bylo zavedení Ceph, rozšířeného open-source úložného řešení, které umožňuje vysoce škálovatelné blokové úložiště.

Proxmox se také vyznačuje svým uživatelským rozhraním založeným na webovém prohlížeči, které umožňuje jednoduché správu a monitoring virtuálních strojů a úložiště. Uživatelé oceňují jeho flexibilitu a širokou škálu konfiguračních možností, které jsou přístupné i bez hlubokých znalostí Linuxu.

Proxmox VE je také známý svou komunitou a širokou podporou. Jako open-source řešení má aktivní komunitu vývojářů a uživatelů, kteří neustále přispívají k jeho vývoji a zlepšování. To zahrnuje pravidelné aktualizace, bezpečnostní opravy a nové funkce, což udržuje Proxmox VE jako relevantní a konkurenceschopnou platformu v oblasti virtualizace.

V současnosti je Proxmox široce používán v mnoha odvětvích, od malých podniků až po velké korporace, díky své schopnosti efektivně spravovat rozsáhlé a komplexní virtualizované prostředí. Jeho kombinace výkonu, flexibility a nízkých nákladů činí z Proxmoxu atraktivní volbu pro mnoho organizací hledajících efektivní virtualizační řešení.
\chapter{virtualizační jádro}
Proxmox VE (Virtual Environment) představuje jeden z předních nástrojů v oblasti virtualizačních technologií, stojící na základech Linuxu a jeho pokročilých funkcí. Jádrem Proxmoxu jsou technologie LXC (Linux Containers) a KVM (Kernel-based Virtual Machine), které jsou integrálními součástmi Linuxového jádra již od verze 2.6. Tyto technologie poskytují základ pro vytváření vysoce výkonného a škálovatelného prostředí, které je ideální pro správu virtualizovaných datových center.

KVM v Proxmoxu umožňuje plnou hardwarovou virtualizaci, díky čemuž mohou být virtuální stroje (VM) spouštěny s vlastním izolovaným operačním systémem, nabízejícím výkon srovnatelný s fyzickými servery. Na druhou stranu, LXC poskytuje lehkou virtualizaci na úrovni kontejnerů, která je vhodná pro aplikace vyžadující rychlé nasazení a efektivní využití zdrojů.

Proxmox VE vyniká také svým intuitivním webovým uživatelským rozhraním, které zjednodušuje správu virtualizovaných prostředí. Správci mohou snadno vytvářet, konfigurovat a spravovat virtuální stroje a kontejnery, stejně jako řídit síťové a úložné zdroje. Tato centralizovaná správa zvyšuje efektivitu a snižuje komplexitu správy datových center.

Další významnou vlastností Proxmoxu je jeho vysoká dostupnost a podpora pro clusterování, což umožňuje vytvářet redundantní a odolné systémy. Proxmox podporuje různé typy úložných řešení, včetně Ceph, NFS, iSCSI, což umožňuje flexibilní a škálovatelné možnosti pro ukládání dat.

Za svou popularitu vděčí Proxmox VE také aktivní a rozsáhlé komunitě uživatelů a vývojářů, kteří neustále přispívají k jeho rozvoji a inovaci. Díky tomu je Proxmox VE neustále na špici technologického pokroku, poskytujíc nejnovější funkce a zabezpečení pro efektivní správu datových center.
%%% LXC
\input{\sectionpath{01}/lxc}

%%% KVM
\input{\sectionpath{01}/kvm}
\chapter{Clusterování}
Clusterování v Proxmoxu VE (Virtual Environment) je jednou z klíčových funkcí, které tuto platformu odlišují jako robustní řešení pro virtualizaci a správu datových center. Clusterování v Proxmoxu umožňuje spojit více fyzických serverů do jednoho logického celku, což vede k zvýšené dostupnosti, lepšímu využití zdrojů a jednodušší správě.

Clusterování v Proxmoxu VE (Virtual Environment) je zásadní funkcí, která transformuje tuto platformu v robustní řešení pro správu a virtualizaci datových center. Tato funkce umožňuje spojení několika fyzických serverů do jednoho logického celku, což zvyšuje dostupnost zdrojů, usnadňuje správu a zlepšuje celkové využití zdrojů.

Nejvýznamnější výhodou clusterování v Proxmoxu je poskytování vysoké dostupnosti. Clusterové řešení automaticky detekuje selhání serveru a přesunuje virtuální stroje na jiný server v clusteru, čímž minimalizuje výpadky služeb. K dosažení tohoto cíle Proxmox využívá quorum mechanismus, který zajistí dostupnost služeb i v případě výpadku některých uzlů.

Proxmox také umožňuje centralizovanou správu všech serverů v clusteru prostřednictvím jednoho webového rozhraní. Tato centralizace umožňuje správcům pohodlně monitorovat a spravovat zdroje, virtuální stroje, kontejnery, síťová nastavení a úložiště pro celý cluster. Všechny změny jsou automaticky replikovány napříč clusterem, což vede k efektivnější správě a nižší pravděpodobnosti chyby v konfiguraci.

Další klíčovou vlastností clusterů v Proxmoxu je podpora pro sdílené úložiště, což umožňuje snadnou migraci virtuálních strojů mezi fyzickými servery. Platforma podporuje různé typy úložišť, včetně lokálních disků, NFS, iSCSI a Ceph, což správcům umožňuje vybrat nejlepší úložnou konfiguraci podle jejich specifických potřeb.

Clusterové řešení v Proxmoxu je navrženo tak, aby bylo škálovatelné a flexibilní. Správci mohou snadno přidávat nebo odebírat servery z clusteru bez výrazného narušení provozu. Tato flexibilita zajišťuje, že datacentra mohou růst a přizpůsobovat se měnícím se požadavkům bez nutnosti provádění rozsáhlých hardwarových nebo softwarových změn.

\pagebreak

Bezpečnost a izolace jsou dalšími důležitými aspekty clusterování v Proxmoxu. Platforma nabízí pokročilé funkce zabezpečení, včetně firewalů, izolace sítě a šifrování, což zajišťuje ochranu dat a služeb v clusteru.

Proxmox, jako open-source řešení, těží z silné komunitní podpory, která přispívá k neustálému vývoji a vylepšování clusterových funkcí. Toto vede k pravidelným aktualizacím a inovacím, udržujícím clusterové řešení Proxmox na špici technologického pokroku.

Clusterování v Proxmoxu VE je nejenom o vysoké dostupnosti a centralizované správě, ale také o efektivním využívání zdrojů a optimalizaci výkonu. Integrace s různými typy úložných řešení, jako jsou Ceph, NFS a iSCSI, umožňuje správcům vytvářet vysoce dostupné a flexibilní úložné systémy, které jsou klíčové pro udržení nepřetržitého provozu služeb. Tato flexibilita v úložných řešeních umožňuje Proxmoxu efektivně reagovat na různé nároky aplikací a zatížení.

Další významnou vlastností clusterování v Proxmoxu je podpora pro live migraci virtuálních strojů a kontejnerů bez výpadku služby. Tato funkce je neocenitelná pro údržbu a upgrade hardwaru, neboť umožňuje přesunutí běžících instancí mezi uzly clusteru bez nutnosti vypínání a restartování, což minimalizuje přerušení služeb a zvyšuje flexibilitu při správě zdrojů.

Proxmox VE nabízí také rozsáhlé možnosti konfigurace a přizpůsobení. Skrze jeho modulární architekturu a rozhraní API mohou správci integraci Proxmoxu s existujícími systémy a procesy automatizovat, což zvyšuje efektivitu a snižuje možnost lidské chyby. Rozhraní API také umožňuje integraci s různými nástroji pro monitoring a správu, což zajišťuje, že správci mají přehled o stavu a výkonu celého clusteru.

Z hlediska bezpečnosti Proxmox VE poskytuje rozsáhlé nástroje pro zabezpečení clusteru a jeho zdrojů. Funkce jako integrovaný firewall, podpora šifrování a izolace sítě zajišťují, že všechny data a komunikace mezi uzly v clusteru jsou chráněny. Tato bezpečnostní opatření jsou zásadní pro ochranu před vnějšími hrozbami a udržení integrity dat.

Komunita okolo Proxmoxu hraje klíčovou roli v jeho vývoji a inovaci. Aktivní uživatelé a vývojáři neustále přispívají k vylepšení platformy, přinášejí nové funkce a udržují Proxmox aktualizovaný vzhledem k nejnovějším technologickým trendům. Tato otevřená komunitní podpora je základem pro udržitelný rozvoj a inovaci Proxmoxu.

\chapter{Uložistě}

Proxmox nabízí širokou škálu možností úložiště, které jsou klíčové pro efektivní správu dat a zdrojů v virtualizovaných prostředích. Tyto možnosti úložiště se pohybují od lokálního úložiště po sofistikovaná síťová řešení, což umožňuje uživatelům Proxmoxu přizpůsobit úložiště podle specifických potřeb a požadavků jejich aplikací a služeb.

Dále také umožňuje použití lokálních disků na hostitelských serverech pro ukládání dat. Toto řešení je vhodné pro menší nasazení nebo pro úložiště, které nevyžaduje vysokou dostupnost. Lokální úložiště poskytuje výhody, jako jsou nízké latence a vysoký výkon, avšak nemá vlastnosti vysoké dostupnosti, které jsou kritické pro větší a náročnější prostředí.

Proxmox podporuje síťové úložiště, jako jsou NFS (Network File System) a iSCSI (Internet Small Computer Systems Interface). Tato úložiště jsou ideální pro prostředí vyžadující vysokou dostupnost a sdílení dat mezi více servery. NFS je často využíván pro sdílení souborů, zatímco iSCSI poskytuje blokové úložiště, což je užitečné pro výkonné databázové servery a aplikace s vysokými nároky na I/O operace.

Ceph je vysoce škálovatelné distribuované úložiště, které se stává čím dál oblíbenějším v prostředí Proxmoxu. Nabízí vysokou dostupnost, vynikající škálovatelnost a vlastnosti samoopravy. Ceph je ideální pro velká a dynamická prostředí, kde je potřeba pružně reagovat na měnící se požadavky na úložný prostor a výkon.

Proxmox VE také podporuje ZFS, pokročilý souborový systém a logický správce svazků. ZFS nabízí integrované funkce, jako jsou snímky, replikace a šifrování. ZFS je obzvláště užitečný pro zálohování a archivaci dat díky jeho schopnosti efektivně spravovat velké objemy dat a poskytovat robustní ochranu proti poškození dat.

Proxmox VE zahrnuje nástroje pro zálohování a obnovu virtuálních strojů a kontejnerů. Tyto nástroje umožňují pravidelně plánovat zálohy, což je zásadní pro ochranu dat a rychlou obnovu v případě selhání systému.

Ve výsledku Proxmox VE poskytuje flexibilní a škálovatelné úložné řešení, které se dá přizpůsobit různým potřebám a požadavkům datových center. Od lokálního úložiště pro rychlé a jednoduché nasazení až po komplexní distribuované systémy, jako je Ceph pro velké a dynamické prostředí, nabízí Proxmox širokou paletu možností pro efektivní správu dat.


\input{\sectionpath{03}/ceph}


\chapter{Zálohování a obnova}

Zálohování lze provést již na platformě Proxmox VE, ale tato metoda není doporučovaná. Pro bezpečnost a dostipnost záloh byla vyvinuta platforma Proxmox Backup Server, která využívá úložistě Ceph pro zaručení dostupnosti záloh. Tato platforma přináší integrované a efektivní řešení pro zálohování virtuálních strojů (VM) a kontejnerů, poskytujíc robustní ochranu dat v prostředích virtualizace.

Proxmox Backup Server je navržen tak, aby efektivně zálohoval VM a kontejnery spravované v Proxmox VE. Umožňuje uživatelům vytvářet plné i inkrementální zálohy, což minimalizuje nároky na úložný prostor a zajišťuje, že zálohy jsou rychlé a efektivní. Inkrementální zálohy zaznamenávají pouze změny od posledního zálohování, což snižuje objem přenesených a uložených dat.

Bezpečnost je klíčovou součástí Proxmox Backup Serveru. Nabízí šifrování na straně klienta, což znamená, že data jsou šifrována před odesláním na zálohovací server a zůstávají šifrována během uložení. Toto zajišťuje vysokou úroveň ochrany dat a je nezbytné pro splnění přísných bezpečnostních požadavků.

Proxmox Backup Server využívá deduplikaci a kompresi dat, aby maximalizoval účinnost úložiště. Deduplikace umožňuje serveru identifikovat a odstranit duplicitní data napříč všemi zálohami, což výrazně snižuje potřebný úložný prostor. Kompresie dále redukuje velikost dat, což zvyšuje efektivitu přenosu dat a úložiště.

Proxmox Backup Server podporuje různé typy úložišť, včetně lokálních disků, NFS a Ceph. To umožňuje uživatelům volit úložiště, které nejlépe vyhovuje jejich specifickým potřebám a rozpočtu.

Proxmox Backup Server je snadno integrovatelný s existujícím Proxmox VE prostředím. Nabízí jednoduché webové rozhraní pro správu zálohovacích úloh, sledování stavu záloh a obnovu dat. Toto intuitivní rozhraní usnadňuje správu záloh i pro uživatele, kteří nejsou odborníky na zálohování.

Obnova dat je klíčovou funkcí Proxmox Backup Serveru. Umožňuje rychlou a snadnou obnovu celých VM, jednotlivých disků nebo souborů. Díky tomu je možné efektivně reagovat na různé scénáře ztráty dat, od jednoduchých chyb uživatele po vážná hardwarová selhání.
\chapter{Migrace a flexibilita}
Zálohování v Proxmox VE (Virtual Environment) je kritickým prvkem pro ochranu dat a virtuálních prostředí v rámci datových center. Tato platforma poskytuje komplexní a efektivní řešení pro zálohování virtuálních strojů (VM) a kontejnerů, což umožňuje organizacím minimalizovat riziko datových ztrát a rychle reagovat na různé scénáře havárií.

Centrální správa záloh je jednou z klíčových vlastností Proxmox VE. Nabízí centrální správu všech záloh z jednoho webového rozhraní, což znamená, že správci systémů mohou pohodlně monitorovat a spravovat všechny zálohy z jediného místa. Toto centralizované řešení zjednodušuje celý proces zálohování. Další významnou funkcí je plánování zálohovacích úkolů. Platforma umožňuje uživatelům plánovat zálohovací úkoly podle jejich potřeb, včetně nastavení pravidelných záloh. Lze nastavit jak plné, tak inkrementální zálohy, což minimalizuje nároky na úložiště a šetří čas. Proxmox VE také využívá deduplikaci a kompresi dat k optimalizaci úložného prostoru a snížení rychlosti přenosu dat. Deduplikace odstraňuje duplicitní data mezi zálohami, zatímco komprese redukuje objem dat, což vede k efektivnějšímu využití úložného prostoru.

Bezpečnost a šifrování dat jsou dalšími klíčovými aspekty zálohovacího procesu v Proxmox VE. Platforma poskytuje šifrování dat během jejich přenosu i uložení, čímž zajišťuje, že citlivá data zůstanou chráněna před neoprávněným přístupem.Rychlá obnova dat je další důležitou funkcí. V případě potřeby lze rychle obnovit celé virtuální stroje, jednotlivé disky nebo soubory ze zálohovaných VM. Tato schopnost umožňuje organizacím rychleji se vrátit do provozu po incidentu.Integrace zálohovacích funkcí s Proxmox VE je také klíčová. Zálohování a správa zálohovacích úkolů jsou plně integrovány přímo do prostředí Proxmox VE, což usnadňuje správu a zvyšuje efektivitu práce.

Nakonec, Proxmox VE jako open-source řešení nabízí výhody spojené s aktivní komunitou uživatelů a vývojářů. Existuje široká podpora a možnost spolupráce na vylepšení platformy, což přináší neustálé inovace a zlepšování.Zálohování v Proxmox VE je tak kritickým prvkem pro zajištění integrity dat a rychlou obnovu v případě nečekaných událostí. Díky funkcím jako šifrování, deduplikace, centralizovaná správa a snadná obnova dat mohou organizace provozovat svá virtuální prostředí s větší jistotou a minimalizovat výpadky služeb.
\chapter{Porovnaní Proxmox a VMware}
Porovnání Proxmox VE a VMware ESXi je klíčové pro pochopení, jak tyto dvě platformy pro virtualizaci nabízí různé možnosti a výhody pro správu datových center.

Proxmox VE je open-source řešení, které poskytuje virtualizaci pomocí KVM a kontejnerizaci prostřednictvím LXC. Je známé svou flexibilitou a širokým spektrem funkcí. Na druhé straně, VMware ESXi je průmyslovým standardem v oblasti virtualizace, nabízející robustní a stabilní platformu s širokou podporou pro různé operační systémy a aplikace.

Co se týče uživatelského rozhraní a správy, Proxmox VE nabízí webové rozhraní, které je intuitivní a snadno se ovládá, a umožňuje správu VM a kontejnerů, clustrování, zálohování a další. VMware ESXi naopak používá vSphere pro správu, což je mocný nástroj, ale může být složitější na naučení a používání, zejména pro nové uživatele.

Ve výkonu a škálovatelnosti je Proxmox VE efektivní a škálovatelný, umožňuje snadnou správu zdrojů a jejich rozšiřování. VMware ESXi je známý vysokým výkonem a škálovatelností, což je důvod, proč je preferován v náročnějších podnikových prostředích.

Podpora a komunita jsou také důležité aspekty. Proxmox VE má silnou a aktivní open-source komunitu, která poskytuje bezplatnou podporu a bohaté zdroje pro uživatele. Na druhé straně, VMware má rozsáhlou profesionální podporu a rozsáhlou uživatelskou základnu, což je klíčové pro podniky vyžadující garantovanou podporu a služby.

Licencování a náklady jsou další důležité faktory. Proxmox VE je bezplatný, s volitelným předplatným pro profesionální podporu a přístup k enterprise úložištím. Naproti tomu, VMware ESXi vyžaduje licenci, což může znamenat významné náklady, zejména ve větších nasazeních.

Proxmox VE nabízí vysokou úroveň přizpůsobitelnosti a množství doplňků, včetně silné podpory pro zálohování a obnovu. VMware poskytuje širokou škálu funkcí a integrací s dalšími produkty VMware, což z něj činí komplexní řešení pro podniková prostředí.

Migrace a flexibilita jsou další klíčové aspekty. Proxmox VE umožňuje snadnou migraci VM a kontejnerů mezi hosty a podporuje živou migraci. VMware ESXi také podporuje živou migraci VM, ale může být více omezený vzhledem k hardwarovým a licenčním požadavkům.

Bezpečnost je také důležitá. Proxmox VE nabízí dobré bezpečnostní funkce, ale jako open-source řešení může vyžadovat více úsilí pro nastavení a udržování bezpečnosti. VMware ESXi je široce považován za velmi bezpečnou platformu s pravidelnými aktualizacemi a siln

ou bezpečnostní architekturou.

Výběr mezi Proxmox VE a VMware ESXi závisí na konkrétních potřebách organizace. Pro malé a střední podniky nebo pro ty, kteří preferují open-source řešení, může být Proxmox VE vhodnější. Pro velké podniky, které vyžadují rozsáhlé funkce, vysokou stabilitu a profesionální podporu, je často preferován VMware ESXi. Obě platformy nabízejí silné řešení pro virtualizaci s různými výhodami a omezeními.
\chapter{Závěr}

Proxmox VE, jako open-source platforma, nabízící rozsáhlou podporu a flexibilitu pro řízení virtualizovaných prostředí. Jeho schopnost integrace s různými úložnými řešeními, včetně Ceph, a podpora pro KVM a LXC, přispívají k jeho adaptabilitě pro různé podnikové potřeby. Zdůrazňuji zde také, že Proxmox VE zvyšuje efektivitu správy datových center a poskytuje bezpečné a škálovatelné řešení pro podniky všech velikostí.

%%% Seznam použité literatury
\nocite{*}
\printbibliography

\openright
\end{document}